% % % % % % % % % % % % % % % % % % % % % % % % % % %
% IS&T Template
% ECE5110 Unit 03 Article
% % % % % % % % % % % % % % % % % % % % % % % % % % %

%%%%%%%%%%%%%%%%%%%%%%%%%%%%%%%%%%
% Document class
%%%%%%%%%%%%%%%%%%%%%%%%%%%%%%%%%%
\documentclass[letterpaper,twocolumn,fleqn]{article}

%%%%%%%%%%%%%%%%%%%%%%%%%%%%%%%%%%
% Packages
%%%%%%%%%%%%%%%%%%%%%%%%%%%%%%%%%%
\usepackage{ist}
\usepackage{amsmath}
\usepackage{amssymb}
\usepackage{siunitx}
\usepackage{physics}
\AtBeginDocument{\RenewCommandCopy\qty\SI}
\ExplSyntaxOn
\msg_redirect_name:nnn { siunitx } { physics-pkg } { none }
\ExplSyntaxOff

\pagestyle{empty}

%%%%%%%%%%%%%%%%%%%%%%%%%%%%%%%%%%
% Title and Authors
%%%%%%%%%%%%%%%%%%%%%%%%%%%%%%%%%%
\title{Numerical Integration and Differentiation of $x^2$ Using Newton--Cotes Rules and Finite Differences}
\author{Juan Ortiz; California Polytechnic Pomona; Pomona, California}
\date{\today}

%%%%%%%%%%%%%%%%%%%%%%%%%%%%%%%%%%
% Begin document
%%%%%%%%%%%%%%%%%%%%%%%%%%%%%%%%%%
\begin{document}
\maketitle
\thispagestyle{empty}

%%%%%%%%%%%%%%%%%%%%%%%%%%%%%%%%%%
% Abstract
%%%%%%%%%%%%%%%%%%%%%%%%%%%%%%%%%%
\begin{abstract}
This article presents a validation exercise for numerical integration and differentiation using the benchmark function $f(x)=x^2$. The integral over $[0,1]$ is approximated with the Midpoint, Trapezoidal, and Simpson rules and compared against the exact value $1/3$. The derivative at $x=1$ is estimated using central, forward, and backward finite differences and compared against the exact slope $2$. The reported Python results show near machine-precision agreement for Simpson's Rule and superior derivative accuracy for the central difference method.
\end{abstract}

%%%%%%%%%%%%%%%%%%%%%%%%%%%%%%%%%%
% Introduction
%%%%%%%%%%%%%%%%%%%%%%%%%%%%%%%%%%
\section{Introduction}
\label{sec:intro}
Numerical methods are essential in electrical engineering when analytic integration or differentiation is difficult or unavailable. A standard way to validate implementation quality is to test methods on a polynomial with known exact results. In this study, the function
\[
f(x)=x^2, \quad x\in[0,1]
\]
is used to evaluate quadrature and finite-difference behavior before extending the same workflow to more complex ECE5110 modeling tasks.

%%%%%%%%%%%%%%%%%%%%%%%%%%%%%%%%%%
% Theory / Method
%%%%%%%%%%%%%%%%%%%%%%%%%%%%%%%%%%
\section{Theory and Method}
The exact reference quantities are
\[
\int_0^1 x^2\,\dd x = \left.\frac{x^3}{3}\right|_{0}^{1}=\frac{1}{3}=0.3333333333333333,
\]
\[
\dv{}{x}(x^2)=2x, \quad \left.\dv{}{x}(x^2)\right|_{x=1}=2.
\]

The numerical integration methods applied are the Midpoint Rule, Trapezoidal Rule, and Simpson's Rule. Numerical differentiation at $x=1$ is evaluated with central, forward, and backward differences. For each method, absolute error is reported relative to the exact analytic value.

%%%%%%%%%%%%%%%%%%%%%%%%%%%%%%%%%%
% Application / Analysis
%%%%%%%%%%%%%%%%%%%%%%%%%%%%%%%%%%
\section{Application and Analysis}
The Python implementation outputs numerical approximations and absolute errors for both tasks. The integration case evaluates the definite integral of $x^2$ on $[0,1]$, while the differentiation case evaluates the slope at $x=1$. This paired test verifies both area-under-curve and local-rate-of-change operators in a single controlled benchmark.

%%%%%%%%%%%%%%%%%%%%%%%%%%%%%%%%%%
% Error Analysis
%%%%%%%%%%%%%%%%%%%%%%%%%%%%%%%%%%
\section{Error Analysis}
For the integration task, Simpson's Rule has the smallest error:
\[
\abs{E_S}=\num{5.551115123125783e-17},
\]
which is effectively at floating-point precision for this case. The Midpoint and Trapezoidal rules also converge closely, with errors on the order of \num{1e-9}.

For numerical differentiation, the central difference error
\[
\abs{E_C}=\num{2.000177801164682e-12}
\]
is lower than both one-sided approximations, consistent with the expected higher-order behavior of the central scheme.

%%%%%%%%%%%%%%%%%%%%%%%%%%%%%%%%%%
% Numerical Comparison / Results
%%%%%%%%%%%%%%%%%%%%%%%%%%%%%%%%%%
\section{Numerical Comparison and Results}
Table~\ref{tab:integration} lists the integration outputs.

\begin{table}[h!]
\caption{Integration results for $\int_0^1 x^2\,\dd x$}
\label{tab:integration}
\centering
{\small
\setlength{\tabcolsep}{3pt}
\sisetup{round-mode=figures,round-precision=8}
\begin{tabular}{|l|c|c|}
\hline
Method & Approx. & Abs. Error \\\hline
Midpoint & $\num{0.3333333325000001}$ & $\num{8.333332357501888e-10}$ \\\hline
Trapezoidal & $\num{0.33333333500000006}$ & $\num{1.6666667490561338e-9}$ \\\hline
Simpson & $\num{0.33333333333333337}$ & $\num{5.551115123125783e-17}$ \\\hline
Exact & $\num{0.3333333333333333}$ & $\num{0}$ \\\hline
\end{tabular}
}
\end{table}

Table~\ref{tab:diff} lists the derivative outputs at $x=1$.

\begin{table}[h!]
\caption{Differentiation results for $f'(1)$ where $f(x)=x^2$}
\label{tab:diff}
\centering
{\small
\setlength{\tabcolsep}{3pt}
\sisetup{round-mode=figures,round-precision=8}
\begin{tabular}{|l|c|c|}
\hline
Method & Approx. & Abs. Error \\\hline
Central diff. & $\num{2.000000000002}$ & $\num{2.000177801164682e-12}$ \\\hline
Forward diff. & $\num{2.000000000035307}$ & $\num{3.530686853991938e-11}$ \\\hline
Backward diff. & $\num{1.9999999999686933}$ & $\num{3.130673498219494e-11}$ \\\hline
Exact & $\num{2.0}$ & $\num{0}$ \\\hline
\end{tabular}
}
\end{table}

%%%%%%%%%%%%%%%%%%%%%%%%%%%%%%%%%%
% Discussion
%%%%%%%%%%%%%%%%%%%%%%%%%%%%%%%%%%
\section{Discussion}
The computed values follow theoretical expectations for smooth polynomial data. Simpson's Rule is most accurate for the integral in this benchmark, while central differencing outperforms forward and backward formulas for slope estimation at $x=1$. This agreement between analysis and implementation supports using the same numerical pipeline for subsequent ECE5110 problems.

%%%%%%%%%%%%%%%%%%%%%%%%%%%%%%%%%%
% Conclusion
%%%%%%%%%%%%%%%%%%%%%%%%%%%%%%%%%%
\section{Conclusion}
Using $f(x)=x^2$ as a validation function, the numerical workflow demonstrates high fidelity against exact analytic results. Simpson's Rule achieves near machine-precision integration accuracy, and central differencing provides the best derivative estimate among tested finite-difference methods. These outcomes establish a reliable baseline for future numerical modeling studies.

%%%%%%%%%%%%%%%%%%%%%%%%%%%%%%%%%%
% References
%%%%%%%%%%%%%%%%%%%%%%%%%%%%%%%%%%
\small
\begin{thebibliography}{9}
\bibitem{burden}
R. L. Burden and J. D. Faires, \emph{Numerical Analysis}, 10th ed., Cengage Learning, 2015.

\bibitem{chapra}
S. C. Chapra and R. P. Canale, \emph{Numerical Methods for Engineers}, 8th ed., McGraw-Hill, 2020.
\end{thebibliography}

\end{document}
